(1) Internal invariance of $I$ is structural: $I$ is a section of the tractor factor and the internal groups act on the internal factor, and $F_i$ is a two-form with values in the internal Lie algebra, with no tractor slot for $\epsilon$ or $I$ to contract. Topologicality: $d\,\tr(F\wedge F) = \tr(DF\wedge F) + \tr(F\wedge DF) = 2\tr(DF\wedge F) = 0$ by the Bianchi identity $DF = 0$, using that two-forms commute under the wedge and the trace is cyclic; locally $\tr(F\wedge F) = d(\mathrm{CS}(A))$, and the variation $\delta\int\tr(F\wedge F) = 2\int\tr(D\delta A\wedge F) = 2\int d\,\tr(\delta A\wedge F) - 2\int\tr(\delta A\wedge DF)$ is a boundary term plus zero: no field equations. (2) The star on two-forms in four dimensions is $*\alpha_{ab} = \tfrac12\sqrt{|g|}\,\epsilon_{abcd}\,g^{ce}g^{df}\alpha_{ef}$; under $g \mapsto \Omega^2 g$ the factor $\sqrt{|g|}$ scales by $\Omega^4$ and the two inverse metrics by $\Omega^{-4}$, so $*$ is invariant: it depends on $[g]$ only, and not on the scale that the frozen point fixes. It does depend on $[g]$: with $g$ degenerate, $g^{-1}$ does not exist. (3) is (2) read at a corank-one wall, where the cometric $g^{ce}$ has a one-dimensional radical and $\sqrt{|g|}$ a half-power pole (Fact \claimref{dark-energy.imported-walls-sphere-theorem}), for every representative of the conformal class alike, so $*$ on two-forms is not a nondegenerate map there; the $\theta$-terms $\theta\int\tr(F\wedge F)$ are star-free, hence cross. The trace identities used in (1) are confirmed by the script on exact instances; they are the only computational content.
